\section{Struttura della tesi}
\label{sec:struttura}

Il \textbf{Capitolo~\ref{cap:Stato dell'arte}} analizza lo stato dell'arte
relativo all'elaborazione serverless sostenibile, esaminando le
principali soluzioni proposte in letteratura per il monitoraggio
energetico del software, con particolare attenzione agli strumenti
capaci di operare a granularità container-level. Vengono inoltre
discusse le tecniche di scheduling energy-aware per ambienti FaaS
e le direzioni emergenti nell'ambito dell'Approximate Computing
applicato all'esecuzione serverless. Il capitolo si conclude con
la presentazione della piattaforma Serverledge, adottata come base
sperimentale per lo sviluppo del sistema proposto.

Il \textbf{Capitolo~\ref{cap:ProgettazioneRealizzazione}} descrive le scelte
progettuali e i meccanismi realizzati. Vengono presentati il modello
di funzione logica con varianti energetiche, l'integrazione di Kepler
per il monitoraggio dei consumi a livello di container, il worker
periodico responsabile dell'aggiornamento adattivo dei profili
energetici e la logica decisionale dello scheduler energy-aware,
con particolare attenzione al ruolo del consenso esplicito
dell'utente nella selezione di varianti approssimate.

Il \textbf{Capitolo~\ref{cap:RisultatiSperimentali}} riporta i risultati
della valutazione sperimentale condotta su un insieme di funzioni
di test rappresentative di differenti profili computazionali.
Vengono analizzati il comportamento dello scheduler al variare
del budget energetico, il risparmio energetico ottenuto tramite
le varianti approssimate e il corrispondente impatto sull'accuratezza
del risultato, quantificando il compromesso introdotto dalle
scelte progettuali adottate.

Il \textbf{Capitolo~\ref{cap:Conclusioni}} raccoglie le conclusioni
del lavoro, discutendo i risultati ottenuti in relazione agli
obiettivi definiti e identificando le principali direzioni di
sviluppo futuro.
.
